\begin{mdframed}
    \textbf{La extensión máxima para esta sección es un párrafo breve (máximo 1/4 de página).}
\end{mdframed}

La experiencia permitió obtener conclusiones sobre a la complejidad teórica de los algoritmos respecto a su implementación práctica.

En el caso del ordenamiento, los cuatro algoritmos evaluados comparten la misma complejidad teó-
rica $\mathcal{O}(n \log n)$, a exepción de \textit{Patiencesort}, pero presentan diferencias significativas en tiempo de ejecución para entradas grandes.
\textit{Sort} resulta ser el más eficiente en la práctica debido a las optimizaciones implementadas en la
STL, seguido de \textit{Quicksort} con mediana de tres. \textit{Mergesort} y \textit{Patiencesort}, a pesar de tener complejidad
asintótica similar, son los más lento para $n \leq 10^7$ por el overhead de memoria que introducen sus implementaciones.

 En el caso de la multiplicación de matrices, \textit{Strassen} demuestra que una menor complejidad asintó-
 tica no garantiza un mejor rendimiento práctico. Las constantes ocultas en la notación Big O afectan en el comportamiento de los algoritmos para los tamaños evaluados.
 
 Todo esto considerando que las implementaciones no son las mejor optimizadas, estos resultados refuerzan la importancia del análisis experimental como complemento al análisis teórico.